% chapters/12/contoh.tex
\section{Dua Contoh Terapan}
\label{sec:contoh-bab-12}

\subsection{Contoh Pertama: Membaca Tangkapan Layar Editor}

\begin{contohbox}
Perhatikan kembali Gambar~\ref{fig:editor-unity}. Tangkapan layar itu
memperlihatkan sebuah \istilah{project} 2D yang baru dibuat, dengan
\istilah{scene} kosong yang memuat tiga benda.
\end{contohbox}

Tabel~\ref{tab:baca-tangkapan} memuat pembacaan setiap bagian tangkapan layar
itu.

\begin{table}[htbp]
\centering
\caption{Pembacaan bagian-bagian tangkapan layar editor}
\label{tab:baca-tangkapan}
{\setlength{\tabcolsep}{4pt}\footnotesize%
\begin{tabularx}{\textwidth}{@{}L{4.0cm} Y@{}}
\toprule
\headrow \textbf{Yang tampak} & \textbf{Keterangan} \\
\midrule
Daftar di kiri atas &
\istilah{Hierarchy}, memuat sebuah scene bernama \istilah{SampleScene} dan tiga
benda di dalamnya. \\
\addlinespace
Ruang abu-abu di tengah, dengan tanda kamera, matahari, dan kubus &
\istilah{Scene}, menampilkan benda-benda itu dari sudut pandang perancang.
Ketiga benda itu adalah kamera, pencahayaan, dan objek primitif berupa kubus. \\
\addlinespace
Panel di kanan dengan tulisan \istilah{Inspector} &
\istilah{Inspector}, menampilkan keterangan mengenai project dan bantuan
pengenalan, sebab belum ada benda yang dipilih. \\
\addlinespace
Daftar berkas di bawah, dengan folder dan berkas &
\istilah{Project window}, memuat berkas project, yaitu scene
\istilah{SampleScene} beserta dua paket bawaan. \\
\bottomrule
\end{tabularx}}
\end{table}

Ketiga benda pada scene itu adalah susunan bawaan yang
disediakan Unity ketika sebuah \istilah{project} 2D dibuat: sebuah kamera, sebuah
pencahayaan, dan sebuah objek primitif. Susunan itulah yang akan Anda susun
sendiri pada praktikum, lalu Anda ubah sesuai keperluan.

\subsection{Contoh Kedua: Menelusuri Urutan Pemanggilan Fungsi}

\begin{contohbox}
Sebuah skrip ditempelkan pada sebuah kubus, lalu scene dijalankan selama dua
detik dengan laju 60 bingkai per detik, kemudian dihentikan. Program
mencetak pesan melalui \code{Debug.Log}.
\end{contohbox}

Urutan pemanggilan fungsinya dapat ditelusuri sebagai berikut.

\begin{enumerate}
  \item \textbf{\code{Awake} dan \code{Start}} dipanggil sekali, sebelum bingkai
        pertama. Pesan \code{Skrip mulai} tercetak satu kali, beserta posisi awal
        benda.
  \item \textbf{\code{Update}} dipanggil 120 kali, yaitu sekali untuk setiap
        bingkai selama dua detik pada 60 bingkai per detik.
  \item \textbf{Pelaporan berkala} tercetak sekali setiap dua detik, yaitu satu
        kali pada akhir penjalanan. Pesannya memuat waktu berjalan, jarak benda
        dari posisi awalnya, dan selang bingkai.
  \item \textbf{\code{OnDisable}} dipanggil sekali ketika penjalanan dihentikan.
        Pesannya memuat lama penjalanan.
\end{enumerate}

Urutan itu seluruhnya ditentukan mesin, bukan oleh kode Anda.
Lama penjalanan diukur oleh skrip Anda sendiri melalui
penjumlahan \code{Time.deltaTime}, sebagaimana kebiasaan yang telah dibangun
sejak Bab~IV: yang tidak diukur tidak dapat diperiksa.

\begin{catatabox}
Perhatikan selang bingkai yang dilaporkan, yaitu sekitar 16,7 milidetik pada 60
bingkai per detik. Angka itu jauh lebih besar daripada selang bingkai pada
program Pygame yang Anda tulis, yang dapat mencapai 5 milidetik ketika tanpa
pembatasan. Bahkan perhatikan bahwa pada Unity, pembatasan laju bingkai
dilakukan mesin, bukan melalui \code{clock.tick()}. Selang yang lebih besar itu
wajar, sebab mesin mengerjakan jauh lebih banyak pekerjaan setiap bingkai.
\end{catatabox}
